\section{Conclusion}
\seclabel{conclusion}

\subsection{Related work}
\subseclabel{related work}

\paragraph{\wLOG in theorem provers}

Given that \WLOG is a common phrase in mathematical proofs, it comes
as no surprise that several major theorem provers feature a similarly
named concept. What is more compelling is that the treatment is far
from uniform, pointing toward a missing uniform foundation. Our
work provides a uniform foundation for symmetric uses of \WLOG{}.

For example, the Rocq library, SSReflect~\cite{ssreflect}, Isabelle/HOL's Archive of Formal Proofs~\cite{unruh:wlog-archive-of-formal-proofs} and the
Lean~4 Mathlib~\cite{lean-mathlib}, define tactics named
$\mathtt{wlog}$, which cut in an additional statement~$\psi$ when proving
$\phi$. Rocq's tactic\footnote{%
  Rocq Reference Manual:
  \tiny\url{https://rocq-prover.org/doc/V9.2.0/refman/proof-engine/ssreflect-proof-language.html\#rocq:tacn.wlog}%
}
%
creates an additional subgoal $(\psi \Rightarrow \phi) \Rightarrow
\phi$.  Isabelle/HOL's tactic similarly adds assumptions \WLOG{}, and
further accounts for generalisation of variables, explicitly
documenting symmetric arguments as the motivation.
Lean's tactic\footnote{%
  Mathlib4 documentation:
  \tiny\url{https://leanprover-community.github.io/mathlib4_docs/Mathlib/Tactic/WLOG.html}
}
%
additionally assumes $\lnot \psi$, adding a subgoal $(\psi \Rightarrow
\phi) \Rightarrow \lnot \psi \Rightarrow \phi$. In all three tactics, the
additional subgoal (e.g., $\psi$ in Rocq and Lean) can be arbitrary and not tied to any
particular symmetry, so it could be worth investigating how many of
its uses correspond to the assumption of canonicity in our framework.

Agda Stdlib defines\footnote{%
  Agda Stdlib:
  \tiny\url{https://github.com/agda/agda-stdlib/blob/c2abce1b8e72ccae49f73791d0996eaf23674357/src/Relation/Binary/Consequences.agda\#L197}
}
%
$\mathtt{wlog} : (\forall (a\, b : A). a \, R \, b \Rightarrow Q \, a
\, b) \Rightarrow (\forall a\, b. Q \, a \, b )$ for any total
relation~$R$ and a symmetric predicate~$Q$. Here, we have a swapping
$\Sym 2$-action on the domain $A \times A$, while $Q$ provides an
appropriate indexed action on the codomain. The
totality of $R$
encodes the canoniser
implicitly, and the assumption exactly
corresponds to our core function.
\citet{dijkstra:wlog} uses a similar setting to criticise \WLOG arguments
for not explicating the underlying symmetry.

A richer collection of symmetries can be found in
\citepos{DBLP:conf/tphol/Harrison09} seminal paper on \WLOG tactics in
HOL. In addition to the tactic discussed in \subsecref{propositional
reasoning}, Harrison develops four additional tactics for \WLOG{}
reasoning in geometry, using a bank of results about invariance and
equivariance of various operations under translations and orthonormal
transformations. Three tactics are immediate special cases of our
principle, while the fourth one is a reduction from a predicate on
points of the form $(\x,\y,0) \in \reals^3$ to a predicate on
$\reals^2$. \citet{schmoetten-palmer-fleuriot:formalising-geometric-axioms-for-minkowski-spacetime-and-without-loss-of-generality-theorems,%
schmoetten-palmer-fleuriot:towards-formalising-schutz-axioms-for-minkowski-spacetime-in-isabelle-hol%
} use Harrison's approach and prove further \WLOG{} lemmata when
formalising Schutz's axioms for Minkowski
spacetime~\cite{schutz:independent-axioms-for-Minkowski-space-time}.

In summary, our work does not account for all
previous principles named $\mathtt{wlog}$. Instead, we offer a crisp
criterion by which one can distinguish symmetric \WLOG{} arguments.

\paragraph{Nominal techniques}

Nominal techniques are a close topic in programming languages, an
established approach to treating bound variables, in which
$\alpha$-equivalence is defined with groups of permutations of atoms acting on terms built from such
atoms~\cite{DBLP:journals/iandc/Pitts03,pitts:nominal-sets-book,DBLP:journals/siglog/Pitts16}.
Our approach borrows from, builds upon, and is analogous to these
nominal techniques. Like nominal techniques, we make pervasive use of, and
rely on, equivariance and related concepts~\cite{pitts:nominal-sets-book}.
Our use of symmetry strengths is similar to \citepos{blanchette-et-al:bindings-as-bounded-natural-functors} use of bounded natural functors.
The proof obligations we
require to establish the existence of recursively defined functions
and their compatibility with the relevant operators is also similar to
the recursors and induction principles in nominal techniques~\cite{%
  gabbay-pitts:a-new-approach-to-abstract-syntax-with-variable-binding,%
  norrish:recursive-function-definition-for-types-with-binders,%
  pitts:alpha-structural-recursion-and-induction,%
  urban-berghofer:a-recursion-combinator-for-nominal-datatypes-implemented-in-isabelle-hol,%
  blanchette-et-al:bindings-as-bounded-natural-functors,%
  popescu:nominal-recursors-as-epi-recursors,%
  popescu-gunter:recursion-principles-for-syntax-with-bindings-and-substitution%
}, and so on. Nominal techniques are typically used
with the group of finitely-supported permutations of names, but not always. The
finite-support requirement is sometimes relaxed~\cite{blanchette-et-al:bindings-as-bounded-natural-functors,
dowek-gabbay-mulligan:permissive-nominal-terms-and-their-unification,
cheney:nominal-logic-programming}.
The names we permute may have additional structure~\cite{bojanczyk-klin-lasota:automata-with-group-actions}.
Nominal metatheory typically employs \WLOG{} reasoning pervasively, cf., e.g., \citepos{cheney:nominal-logic-programming} thesis.

Unlike nominal techniques, our recursion principles do not solely
account for binders and $\alpha$-equivalence. They do require the
additional structure of canonisers. Whereas in nominal techniques one
stays equivariant throughout, in symmetric programming there is no
such requirement and outside the definition of a \WLOG{} construction,
code or reasoning may violate equivariance. We want to revisit this
comparison after we develop symmetric programming and reasoning
further. An even more interesting direction is to combine symmetric
programming and nominal techniques, pervasively maintaining nominal
equivariance, but potentially violating equivariance for non-nominal
symmetries.  Mechanically representing and reasoning about nominal data
is a wide area of research~\cite{%
  poplmark1,%
  poplmark2,%
  schafer-et-al:autosubst,%
  aydemir-et-al:lngen,%
  urban:nominal-isabelle/hol,%
  forster-stark:coq-a-la-carte,%
  kathrin-stark:thesis,%
  stark-et-al:autosubst-2,%
  pientka:beluga,%
  gacek:thesis,%
  gacek-et-al:abella-journal,%
  sewell-et-al:ott,%
  allais-et-al:type-and-scope-safe-syntax-journal,%
  gabbay-pitts:a-new-approach-to-abstract-syntax-with-variable-binding,%
  chlipala:parametric-hoas,%
  chargueraud:locally-nameless,%
  weirich-et-al:binders-unbound,%
  poulsen-et-al:definitional-interpreters,%
  cheney:towards-names,%
  ghani-et-al:cyclic-structures,%
  hamana:initial-semantics-for-cyclic-sharing,%
  polonowski:automatically-generated-syntaxes,%
  keuchel-et-al:needle-and-knot,%
  keuchel-jeuring:generic-conversions,%
  keuchel:msc-thesis,%
  benton-et-al:strongly-typed-terms,%
  lee-et-al:gmeta,%
  erdi:well-scoped-syntax,%
  copello-et-al:variable-convention,%
  morris-et-al:regular-tree-types,%
  hofmann:semantical-analysis-of-ho-abstract-syntax,%
  pfenning-elliott:hoas,%
  altenkirch-reus,%
  blanchette-et-al:bindings-as-bounded-natural-functors}.
While there are first steps towards nominal
type
theories~\cite{aydemir-bohannon-weirich:nominal-reasoning-techniques-in-coq,pitts-matthiesen-derikx:a-dep-type-theory-with-abstractable-names},
we still lack an implementation of a nominal dependently-typed system.


\paragraph{Relational parametricity.}
If we view each symmetry $g \in \G$ as a relation $a \equiv_g b$ given
by $g\act a = b$, then equivariance becomes relational preservation:
$a \equiv_g b \implies f\,a \equiv_g f\,b$. This perspective suggests
some connections with logical
relations~\cite{plotkin:lambda-definability-and-logical-relations} and
with relational
parametricity~\cite{reynolds:types-abstraciton-and-parametric-polymorphism}.
A possible connection is the work of \citet{DBLP:conf/popl/Atkey14} on
using parametricity and Noether's theorem to derive invariants of
physical systems. A shared theme is identifying the group actions that
describe a system's symmetries.  However, applying Noether's theorem
requires a notion of differentiation, which is not necessarily present in our
general setting. One direction for future work is to investigate which
of our action combinators can be translated to smooth-manifold
settings, and whether they can serve as compositional operators for
describing physical systems. It may be fruitful to explore other
connections. For example, there might be some proofs for free
theorems~\cite{wadler:theorems-for-free} that appeal to parametricity
through relations defined via symmetry. There might be symmetric
invariants for restricted programming languages (e.g., with nominal
features) provable via logical relations.

\paragraph{Symmetry in programming languages and type theory}
A more recent development on symmetries in programming is a Hoare
logic for symmetry properties by~\citet{10.1145/3763080}. It features
triples that relate the symmetries of preconditions to those of
postconditions. In particular, a triple of the form $(\G,\A) C
(\H,\B)_\phi$ uses a $\G$-action $\A$ to capture symmetries of input
states and an $\H$-action $\B$ to capture symmetries of output states
after executing the program $C$. The program $C$ is required to be
equivariant with respect to a group homomorphism $\phi : \G \to \H$.
More precisely:
\(
  \forall g \in \G, a \in \carrier\A. C(g \act[\A] a) = \phi(g) \act[\B] C(a)
\).
Beyond general action constructions such as (non-dependent) products
and pre-/post-composition, the paper develops a forward-reasoning
system that simultaneously derives an action on the output and proves
the equivariance of $C$. As we move from mathematical foundations
toward more practical symmetric programming, we
may adapt these rules to ease the construction of the group
actions required by the \WLOG construction.

\paragraph{Symmetry elsewhere in computer science}

Computer science is vast, and group theory is pervasive, so there are
many points of contact between the two. Rather than attempting a
comprehensive review, we highlight a few representative areas and
point to recent surveys for further reading.

In computer algebra and computational group theory, groups are both
the object of study and a computational resource, shaping algorithm
design and complexity analysis in areas such as permutation-group
algorithms and constructive
recognition~\citep{Seress_2003,Holt2005HandbookOC}. The
concept of \emph{canonical labelling} from
nauty/Traces~\cite{MCKAY201494}, which finds the vertex permutation
mapping a graph into its representative element w.r.t. graph
isomorphism, exactly matches our canonisers. It is worth investigating
how many algorithms building on it fit into our \WLOG pattern.

In constraint programming and satisfiability, symmetries have long
been a central concern. A large body of work studies symmetry-breaking
techniques, which add constraints that rule out symmetric solutions
while preserving
satisfiability~\citep{DBLP:reference/fai/GentPP06,DBLP:conf/aaai/Walsh12,janota2025breakingsymmetriesquantifiedgraph}.
In machine learning, symmetry appears both as an
architectural prior (equivariance/invariance) and as an analysis
tool. In computer vision, it supports tasks such as
recognition, reconstruction, and compression of geometric
data~\citep{basheer2025symmetry,zhao2025symmetryneuralnetworkparameter,cornish2025stochasticneuralnetworksymmetrisation}.

Similarly, model checking techniques and tools exploit symmetry to
mitigate the state-space explosion problem~\cite{%
  clarke-jha:symmetry-and-induction-model-checking,%
  donaldson:thesis,%
  clarke2018handbook,%
  miller-donaldson-calder:symmetry-in-temporal-logic-model-checking,%
  sistla-godefroid:symmetry-and-reduced-symmetry-in-model-checking%
}. Model checkers typically try to compute the relevant symmetries~\citep[e.g.][]{
  wahl-d'silva:a-lazy-approach-to-symmetry-reduction,%
  chu-jaffar:a-complete-method-for-symmetry-reduction-in-safety-verification,%
  leuschel-et-al:symmetry-reduction-for-B-by-permutation-flooding,%
  wahl-damm-hermanns:adaptive-symmetry-reduction,%
  bosnacki-dams-holenderski:symmetric-spin%
},
and in the general case this computational task is equivalent to the
graph isomorphism problem~\cite{jha:thesis}, whose precise complexity
class is still open.

\subsection{Future work}


On the foundational side, it may be interesting to explore richer data type constructions, such
as: co-data and co-recursion via terminal co-algebras; and universes.
The theory of distributive laws also applies to $F$-coalgebras,
suggesting there is a similar theory of coinductive
data-types. Universes offer various directions, depending on the
symmetries one may want to equip the inhabitants of the universe.
One option is to equip a universe of types with the trivial action,
as we did with the universe of propositions~(\subsecref{propositional reasoning}).

A richer direction involves a universe $\mathcal U_G$ of $\G$-actions, for some or
each group $\G$.  Here each value of the universe is a code for the
underlying carrier and the $\G$-action. This universe supports the
various type constructors using the usual type-interpretation in the
category of $\G$-actions.  While we can equip this universe with a
trivial $\G$-action, we can also equip it with other actions.

For example, given a symmetry $g \in \underline G$ and an inhabitant  $A \in \mathcal U_G$ with
$\sem{A} = (\underline A, \circledast)$, define $g
\circledast_{\mathcal U_G} A$ as the action with
the carrier unchanged $\underline{g\circledast A} \definedby \underline A$
which acts through conjugation $h \circledast_{g \circledast A} x \definedby
(g^{-1} \ast h \ast g) \circledast x$.  These definitions give
a map $\carrier\G \times \mathcal U_G \to \mathcal U_G$, and moreover this
map is a $G$-action over $\mathcal U_G$.  So in this case there is
some room to explore symmetries on a universe of $G$-actions and
potential to employ symmetric programming/reasoning here. In another direction the
equivariant initial algebra theorem may internalise to an equivariant
function on codes.

In another kind of universe, each inhabitant describes a group $\G$
and a code in a universe $\mathcal U_{[\![G]\!]}$ of $\G$-actions. One
can then similarly look for useful actions and symmetric
programming/reasoning there. Another direction in this area would be
to explore codes for homomorphisms between codes for $G$-actions and
actions and corresponding transport functions between their decoded
$G$-actions and actions. This direction may open up
questions regarding coherence of the resulting type theory.

That said, our development does not need a universe, it is merely
compatible with these various choices, and we plan to explore them in
future work.

On the practical side, we want to refine the developed modular mathematical
abstractions into programming constructs,
guided by more involved symmetric programming and reasoning
examples. In the first instance, we want to complete the full set of
operations on self-balancing binary search
trees~\cite{mcbride:neighbours-in-order}.  Another promising example
is \citepos{escardo:signed-digit-numerals} representation of exact
reals using infinite streams of signed-digit numerals. Operations on
these streams and proofs about these operations contain multiple
symmetric cases, including swapping the roles of symmetric variables,
and reflecting the digits around $0$.  Computable representations of
the constructive Dedekind reals~\cite{kaucher:interval-analysis-in-the-extended-interval-space-IR,bauer-taylor:dedekind-reals-in-abstract-stone-duality,bauer:efficient-computation-with-dedekind-reals} employ symmetric operations on
the cut's lower and upper sets, with symmetric programming potential.

It may be worth exploring
symmetric programming and reasoning for model checking. Since
symmetric programming exposes symmetric program paths explicitly, it
may alleviate the need to automatically detect symmetry when checking
symmetric programs. A more interesting direction is whether a
symmetry-detecting model checker can suggest refactoring opportunities
that use symmetry more explicitly. Another direction involves using
symmetric programming when implementing model checkers or constraint
solvers.

In addition to being interesting in their own right, such examples should guide
us in developing abstractions for performance engineering, as well as on how to most
ergonomically present groups, actions, canonisers, core functions, and
equivariance to the programmer. For example, immediate canonisers
(see~\propositionref{immediate canonical is kernel}) could admit an
interface similar to \emph{lazy
constructors}~\cite{DBLP:journals/pacmpl/LorenzenLSL25}, where we
extend a data type definition with non-canonical constructors together
with their non-trivial canonising symmetries, but list only canonical
cases in pattern matching, since the non-canonical ones are covered by
the \WLOG construction.

We already explored the landscape with two experiments
in Idris and Lean. Having a clean mathematical resolution
of issues that surfaced, such as termination and efficiency, we can
now return to both the programming language and the
proof assistant: not to increase confidence in the
results' correctness, but to develop a library of reusable
results and lay a foundation for future implementation and mechanisation.
